\documentclass{article}
\usepackage[T1]{fontenc}
\usepackage{lmodern}
\usepackage{graphicx}
\usepackage{amsmath,amssymb}
\usepackage[a4paper,margin=2cm]{geometry}
\usepackage[absolute,overlay]{textpos}
\setlength{\TPHorizModule}{1mm}
\setlength{\TPVertModule}{1mm}
\usepackage[indonesian]{babel}
\usepackage{hyperref}
\setlength{\emergencystretch}{2em}
% unit: Halaman 35

\begin{document}

% === Halaman 35 (python/native) ===

• CFB s-bit mengenkripsi plainteks berukuran s bit setiap kalinya, s  b  

(b = ukuran blok dalam satuan bit).

• Dengan kata lain, CFB s-bit memperlakukan cipher blok sama seperti 

pada cipher alir.

• Mode CFB membutuhkan sebuah antrian (shift register) yang 

berukuran sama dengan ukuran blok masukan (b bit) . 

• Tinjau mode CFB 8-bit yang bekerja pada blok berukuran 64-bit  

pada gambar berikut (s = 8, b = 64):

35

\end{document}
